\BeleskeID{03-s004}
% element: 03-s004:e01
% element: 03-s004:e03
% element: 03-s004:d01
Zamena troulaznog I kola zahteva još dva dvoulazna I kola: prvo formira \(BC\), a drugo \(A(BC)\).
Sa početna dva I kola to daje četiri dvoulazna I kola. Ako pakovanje sadrži sva četiri,
dovoljni su čip sa I2 kolima i čip sa invertorima. Ušteda zato zavisi od broja raspoloživih kola u pakovanju.
Putanje od \(B\) i \(C\) do izlaza prolaze kroz dva gejta, dok putanja od \(A\) prolazi samo kroz završni gejt.
Pri izboru realizacije treba proveriti kašnjenje najduže putanje.

% element: 03-s004:e04
Kod I kola spajanje ulaza daje \(XX=X\), pa ono ne zamenjuje invertor.
Invertorska veza NI kola može biti korisna kada takvo kolo već imamo raspoloživo u čipu.
U primeru sa I kolima invertor je i dalje zasebno potreban.

% element: 03-s004:e02
Neiskorišćeni ulazi obuhvataju slobodne priključke korišćenih kola i ulaze potpuno neiskorišćenih delova čipa.
Odgovarajući nivoi određuju se prema vrsti kola i proizvođačkom opisu komponente;
ulazi se ne ostavljaju nepovezani.
